\documentclass[11pt]{article}

\usepackage[margin=1in]{geometry}
\usepackage[T1]{fontenc}
\usepackage[utf8]{inputenc}
\usepackage{lmodern}
\usepackage{amsmath,amssymb,amsthm,mathtools}
\usepackage{microtype}
\usepackage{hyperref}
\usepackage{enumitem}

\hypersetup{
  colorlinks=true,
  linkcolor=blue,
  citecolor=blue,
  urlcolor=blue,
  pdftitle={A Dyadic-to-Integer Beurling Reduction for a Reciprocal Nyman--Beurling Refinement},
  pdfauthor={Karl Schultze}
}

\newtheorem{theorem}{Theorem}
\newtheorem{lemma}{Lemma}
\newtheorem{proposition}{Proposition}
\newtheorem{remark}{Remark}

\newcommand{\Hcal}{\mathcal H}
\newcommand{\Ucal}{\mathcal U}
\newcommand{\Rcal}{\mathcal R}
\newcommand{\dist}{\operatorname{dist}}
\newcommand{\lcm}{\operatorname{lcm}}
\newcommand{\Span}{\operatorname{span}}
\newcommand{\one}{\mathbf 1}

\title{A Dyadic-to-Integer Beurling Reduction\\for a Reciprocal Nyman--Beurling Refinement}
\author{Karl Schultze\\\small Preliminary research note, v0.1}
\date{June 15, 2026}

\begin{document}
\maketitle

\begin{center}
\fbox{\begin{minipage}{0.92\textwidth}
\textbf{Claim status.} This note is \emph{not} a proof of the Riemann Hypothesis. It proves a structural reduction from a dyadic reciprocal no-stalling problem to an integer Beurling/Nyman--Báez-Duarte approximation distance. The remaining input, that this integer distance tends to zero, is the known Riemann-Hypothesis-level Nyman--Beurling/Báez-Duarte obstruction.
\end{minipage}}
\end{center}

\begin{abstract}
We study a dyadic Nyman--Beurling approximation scheme in the Hilbert space
\[
\Hcal=L^2((1,\infty),y^{-2}\,dy)
\]
with Beurling atoms
\[
F_\alpha(y)=\alpha\lfloor y\rfloor-\lfloor \alpha y\rfloor.
\]
For a denominator $Q$, let
\[
V_Q=\Span\{F_{a/Q}:1\le a\le Q-1\},\qquad r_Q=(I-P_Q)1,\qquad x_Q=\|r_Q\|^2.
\]
We introduce the reciprocal subspace
\[
\Rcal_Q=\Span\{F_{1/n}:2\le n\le Q\}.
\]
Since $F_{1/n}\in V_{\Lambda_Q}$ for $\Lambda_Q=\lcm(1,\dots,Q)$, the reciprocal subspace gives an admissible refinement of $V_Q$. We prove a projection dichotomy: if
\[
d_Q=\dist_{\Hcal}(1,\Rcal_Q),
\]
then either $x_Q\le 4d_Q^2$, or
\[
x_Q-x_{\Lambda_Q}\ge \frac14 x_Q^2.
\]
Consequently, if $d_Q\to0$, then $x_Q\to0$ along the refinement chain $Q\mapsto \Lambda_Q$, and the Nyman--Beurling criterion gives RH. Under the isometry $x=1/y$, the atoms $F_{1/n}$ become the pole-cancelled integer Beurling functions
\[
F_{1/n}(1/x)=\left\{\frac1{nx}\right\}-\frac1n\left\{\frac1x\right\}.
\]
Thus the final obstruction is the integer Beurling/Nyman--Báez-Duarte approximation distance, not a new proof of RH.
\end{abstract}

\section{Setup}

Let
\[
\Hcal=L^2((1,\infty),y^{-2}\,dy)
\]
with inner product
\[
\langle f,g\rangle=\int_1^\infty f(y)\overline{g(y)}\,\frac{dy}{y^2}.
\]
For $0<\alpha<1$, define
\[
F_\alpha(y)=\alpha\lfloor y\rfloor-\lfloor \alpha y\rfloor.
\]
For a positive integer $Q$, define
\[
V_Q=\Span\{F_{a/Q}:1\le a\le Q-1\}.
\]
Let $P_Q$ be orthogonal projection onto $V_Q$, and set
\[
r_Q=(I-P_Q)1,\qquad x_Q=\|r_Q\|^2.
\]
The Nyman--Beurling criterion states, in one of its equivalent forms, that suitable approximation of the constant/indicator function by Beurling functions is equivalent to RH; the integer-dilation strengthening is due to Báez-Duarte \cite{baezduarte2003,baezduarte_arxiv}.

\section{The reciprocal subspace}

Define
\[
\Rcal_Q=\Span\{F_{1/n}:2\le n\le Q\}.
\]
Let
\[
\Lambda_Q=\lcm(1,2,\dots,Q).
\]
Since $1/n=(\Lambda_Q/n)/\Lambda_Q$, we have
\[
F_{1/n}\in V_{\Lambda_Q},\qquad 2\le n\le Q.
\]
Therefore
\[
V_Q+\Rcal_Q\subseteq V_{\Lambda_Q}.
\]
This observation removes the need for any rounding or binning error: the reciprocal atoms themselves are valid directions in the larger denominator space.

\section{Unit-step visibility}

Let $\Ucal\subset \Hcal$ be the closed subspace of functions constant on each unit interval $[k,k+1)$, $k=1,2,\dots$.

For $k\le y<k+1$, we have
\[
F_{1/n}(y)=\frac{k}{n}-\left\lfloor\frac{k}{n}\right\rfloor=\left\{\frac{k}{n}\right\}.
\]
Thus $F_{1/n}\in\Ucal$ for every integer $n\ge2$.

\begin{lemma}[Unit-step visibility]
Let $s_Q=P_{\Ucal}r_Q$. Then
\[
\|s_Q\|^2\ge x_Q^2.
\]
\end{lemma}

\begin{proof}
Since $1=P_Q1+r_Q$ and $r_Q\perp V_Q$, we have
\[
\langle r_Q,1\rangle=\langle r_Q,P_Q1+r_Q\rangle=x_Q.
\]
Since $1\in\Ucal$, it follows that
\[
\langle s_Q,1\rangle=\langle P_{\Ucal}r_Q,1\rangle=\langle r_Q,1\rangle=x_Q.
\]
By Cauchy--Schwarz,
\[
x_Q^2=|\langle s_Q,1\rangle|^2\le \|s_Q\|^2\|1\|^2.
\]
But $\|1\|^2=\int_1^\infty y^{-2}\,dy=1$. Hence $\|s_Q\|^2\ge x_Q^2$.
\end{proof}

\section{The projection dichotomy}

Define
\[
d_Q=\dist_{\Hcal}(1,\Rcal_Q).
\]
Let $\psi_Q=P_{\Rcal_Q}1$. Then
\[
\|1-\psi_Q\|=d_Q,\qquad \|\psi_Q\|\le 1.
\]
Since $\Rcal_Q\subset \Ucal$, we have
\[
\langle r_Q,\psi_Q\rangle=\langle s_Q,\psi_Q\rangle.
\]
Furthermore,
\[
\langle s_Q,\psi_Q\rangle=\langle s_Q,1\rangle-\langle s_Q,1-\psi_Q\rangle.
\]
Thus
\[
|\langle r_Q,\psi_Q\rangle|\ge x_Q-\|s_Q\|d_Q.
\]
Since $\|s_Q\|\le\|r_Q\|=\sqrt{x_Q}$,
\[
|\langle r_Q,\psi_Q\rangle|\ge x_Q-\sqrt{x_Q}\,d_Q.
\]

\begin{lemma}[Projection dichotomy]
For every $Q$, either
\[
x_Q\le 4d_Q^2,
\]
or
\[
x_Q-x_{\Lambda_Q}\ge \frac14 x_Q^2.
\]
\end{lemma}

\begin{proof}
If $d_Q>\frac12\sqrt{x_Q}$, then $x_Q<4d_Q^2$.

Otherwise $d_Q\le \frac12\sqrt{x_Q}$, and therefore
\[
|\langle r_Q,\psi_Q\rangle|\ge \frac12 x_Q.
\]
Let
\[
w_Q=(I-P_Q)\psi_Q.
\]
Since $\psi_Q\in\Rcal_Q\subseteq V_{\Lambda_Q}$ and $V_Q\subseteq V_{\Lambda_Q}$, we have
\[
w_Q\in V_{\Lambda_Q}\ominus V_Q.
\]
Also
\[
\langle r_Q,w_Q\rangle=\langle r_Q,\psi_Q\rangle
\]
because $r_Q\perp V_Q$, and
\[
\|w_Q\|\le\|\psi_Q\|\le1.
\]
The projection gain from $V_Q$ to $V_{\Lambda_Q}$ is at least the squared projection onto the direction $w_Q$, hence
\[
x_Q-x_{\Lambda_Q}\ge \frac{|\langle r_Q,w_Q\rangle|^2}{\|w_Q\|^2}\ge \frac14x_Q^2.
\]
\end{proof}

\begin{theorem}[Dyadic-to-integer Beurling reduction]
If $d_Q\to0$, then $x_Q\to0$ along the refinement chain $Q_{j+1}=\Lambda_{Q_j}$. Consequently RH follows by the Nyman--Beurling criterion.
\end{theorem}

\begin{proof}
Apply the projection dichotomy along $Q_{j+1}=\Lambda_{Q_j}$. If the first alternative occurs infinitely often, then $x_{Q_j}\le 4d_{Q_j}^2\to0$ along a subsequence. Since the distances $x_Q$ are nonincreasing under refinement, this forces convergence along the chain. If the second alternative occurs eventually, then
\[
x_{Q_{j+1}}\le x_{Q_j}-\frac14x_{Q_j}^2,
\]
which also implies $x_{Q_j}\to0$. Hence $x_Q\to0$ along the chain, and the Nyman--Beurling criterion gives RH.
\end{proof}

\section{Identification with the integer Beurling distance}

Set $x=1/y$. Then
\[
y^{-2}\,dy=-dx,
\]
so $\Hcal$ is isometric to $L^2(0,1)$.

Let $\rho(t)=\{t\}=t-\lfloor t\rfloor$. Then
\begin{align*}
F_{1/n}(1/x)
&=\frac1n\left\lfloor\frac1x\right\rfloor-\left\lfloor\frac1{nx}\right\rfloor\\
&=\frac1n\left(\frac1x-\rho(1/x)\right)-\left(\frac1{nx}-\rho(1/(nx))\right)\\
&=\rho\left(\frac1{nx}\right)-\frac1n\rho\left(\frac1x\right).
\end{align*}
Therefore $d_Q$ is the distance from $1$ to the span of the integer Beurling functions
\[
\rho\left(\frac1{nx}\right)-\frac1n\rho\left(\frac1x\right),\qquad 2\le n\le Q.
\]
This is the integer Nyman--Báez-Duarte endpoint, written in the pole-cancelled Beurling-function coordinates naturally associated with the atoms $F_{1/n}$.

\section{Mellin formulation}

The Mellin transform identity is
\[
\mathcal M F_\alpha(s)=\frac{\zeta(s)}s(\alpha-\alpha^s).
\]
Thus
\[
\mathcal M F_{1/n}(s)=\frac{\zeta(s)}s\left(\frac1n-\frac1{n^s}\right).
\]
On the Plancherel line $s=1/2+it$, the endpoint distance can be written as
\[
d_Q^2=\inf_{c_2,\dots,c_Q}\frac1{2\pi}\int_{-\infty}^{\infty}\left|1-\zeta\left(\frac12+it\right)\sum_{n=2}^Q c_n\left(\frac1n-\frac1{n^{1/2+it}}\right)\right|^2\frac{dt}{1/4+t^2}.
\]
The remaining approximation problem is therefore a critical-line inverse-zeta approximation problem. Proving $d_Q\to0$ is not supplied by the present note; it is the RH-equivalent Nyman--Beurling/Báez-Duarte obstruction.

\appendix
\section{Exploratory residue-operator symbol}

This appendix records a formal computation that helps explain why the final obstruction has the Riemann zeta function as its symbol. It is not used in the proof of the reduction theorem.

The orthogonal-complement equations for the unit-step sawtooth functions can be written in terms of a sequence $c=(c_k)$ as
\[
\sum_{k\ge1}c_k(k\bmod n)=0,\qquad n\ge2.
\]
Comparing the equations for $n$ and $n+1$ gives the triangular system
\[
nc_n=\sum_{k>n}\left((k\bmod n)-(k\bmod(n+1))\right)c_k.
\]
Writing $d_k=kc_k$ gives
\[
d_n=\sum_{k>n}K_{n,k}d_k,
\]
where
\[
K_{n,k}=\frac{(k\bmod n)-(k\bmod(n+1))}{k}.
\]
Testing formally on power modes $d_k=k^{-s}$ gives the asymptotic symbol
\[
\kappa(s)=1-\left(1-\frac1s\right)\zeta(s).
\]
Formal fixed modes satisfy $\kappa(s)=1$, hence $\zeta(s)=0$ away from the special point $s=1$. The $\ell^2$ threshold for $k^{-s}$ is $\Re(s)>1/2$. This explains why the operator-theoretic obstruction reflects the RH zero-free half-plane. A rigorous use of this symbol in the critical strip would require smoothed Mellin-localized approximate eigenvector estimates.


\section*{Author's note}

I am Karl Schultze, writing as an independent mathematical explorer outside a formal academic position. I work in a practical trade and have been rebuilding my mathematical foundations through self-study, with a long-term hope of understanding number theory more deeply.

The Riemann Hypothesis has been a deeply personal problem for me for many years. Earlier in my life, I was drawn to it in a way that was not yet healthy or grounded. Returning to it now has been different: the goal has been to follow the mathematics honestly, distinguish proof from intuition, and understand exactly where my ideas lead.

This note records the endpoint of that attempt. The dyadic reciprocal idea was carried as far as I could take it, and it led to the known B\'aez--Duarte/Nyman--Beurling boundary rather than to an independent proof of RH. That is still meaningful to me, because it gives a clear and honest stopping point.

I have my closure. I tried to the best of my ability to study the Riemann Hypothesis and look for a solution, and I am ready for someone else to take this problem on.

\section*{Claim status and transparency}

This note does not claim to prove RH. It proves a conditional structural reduction to an integer Beurling/Nyman--Báez-Duarte approximation distance. The remaining assertion $d_Q\to0$ is the RH-level approximation problem. This draft was prepared with AI assistance and should be independently verified before citation or formal submission.

\begin{thebibliography}{9}

\bibitem{baezduarte2003}
L. Báez-Duarte, \emph{A strengthening of the Nyman--Beurling criterion for the Riemann hypothesis}, Rendiconti di Matematica e delle sue Applicazioni, Serie IX, 14 (2003), no. 1, 5--11.

\bibitem{baezduarte_arxiv}
L. Báez-Duarte, \emph{A strengthening of the Nyman--Beurling criterion for the Riemann Hypothesis}, arXiv:math/0202141.

\bibitem{balazard_deroton}
M. Balazard and A. de Roton, \emph{Sur un critère de Báez-Duarte pour l'hypothèse de Riemann}, arXiv:0812.1689.

\end{thebibliography}

\end{document}
